%% Figure: overview of SLDF. Four methods (top) feed one shared core (bottom).
\begin{figure}[t]
\centering
\begin{tikzpicture}[font=\footnotesize, node distance=4mm,
  m/.style={draw, rounded corners, align=center, minimum height=11mm, text width=27mm},
  c/.style={draw, rounded corners, align=center, minimum height=9mm, text width=24mm, fill=black!6},
  >={Stealth[length=2mm]}]
\node[m] (a) {\textbf{A} Classic fuzzing\\SQLsmith, sweeps,\\issue replay};
\node[m, right=of a] (b) {\textbf{B} Feature fuzzers\\triggers, LATERAL,\\delim-join rewrite};
\node[m, right=of b] (cc) {\textbf{C} Hints + LLM\\generators: new code,\\feature cards, dossiers};
\node[m, right=of cc] (d) {\textbf{D} Rule-targeted loop\\rule cards, hypothesis,\\LLM generator};
\node[c, below=10mm of a] (j) {\circled{1} Engine judge\\validity, reach, plans};
\node[c, right=of j] (o) {\circled{2} Oracles\\switch on/off, 5 configs};
\node[c, right=of o] (t) {\circled{3} Triage\\verify, registry,\\minimize, main, search};
\node[c, right=of t] (h) {\circled{4} Person reads,\\files issue / PR};
\draw[->] (j) -- (o); \draw[->] (o) -- (t); \draw[->] (t) -- (h);
\foreach \x in {a,b,cc,d} { \draw[->] (\x.south) -- ++(0,-4mm) -| (o.north); }
\draw[->, dashed] (j.north) -- ++(0,3mm) -| node[pos=0.25, above, font=\scriptsize]{engine feedback (repair rounds)} (d.south);
\end{tikzpicture}
\caption{\tool: four fuzzing methods run side by side and share one core. Generators written by an LLM (Methods C and D)
pass the engine judge \circled{1} first, which also returns measured feedback for the next round of Method D. Every hit
goes through the same oracles and triage, and a person decides what is filed.}
\label{fig:overview}
\end{figure}
